\documentclass[slideColor]{prosper}
\usepackage{amsmath,amsfonts,amsthm}
%\usepackage{epsfig}
%\include{Qcircuit}


\newcommand{\bbC}{{\mathbb{C}}}
\newcommand{\C}{{\mathbb{C}}}
\newcommand{\F}{{\mathbb{F}}}
\newcommand{\N}{{\mathbb{N}}}
\newcommand{\Q}{{\mathbb{Q}}}
\renewcommand{\R}{{\mathbb{R}}}
\newcommand{\Z}{{\mathbb{Z}}}

\newcommand{\cA}{{\mathcal{A}}}
\newcommand{\cD}{{\mathcal{D}}}
\newcommand{\cH}{{\mathcal{H}}}
\newcommand{\cN}{{\mathcal{N}}}
\newcommand{\cP}{{\mathcal{P}}}
\newcommand{\cQ}{{\mathcal{Q}}}
\newcommand{\cS}{{\mathcal{S}}}
\newcommand{\cU}{{\mathcal{U}}}
\newcommand{\cV}{{\mathcal{V}}}

\newcommand{\tr}{\mathop{\mathrm{tr}}}
\newcommand{\rank}{\mathop{\mathrm{rank}}}
\newcommand{\poly}{\mathop{\mathrm{poly}}}
\newcommand{\E}{\mathop{\mbox{$\mathbf{E}$}}}
\newcommand{\schur}{\mathop{\mathrm{Schur}}}
\newcommand{\schursmall}[2]{{\mathrm{S}}_{#1,#2}}
\newcommand{\planch}{\mathop{\mathrm{Planch}}}
\newcommand{\planchsmall}[1]{{\mathrm{P}}_{#1}}
\renewcommand{\d}{{\mathrm{d}}}

\newcommand{\HSP}{\textsc{hsp}\xspace}

\renewcommand{\>}{\rangle}
\newcommand{\<}{\langle}
\newcommand{\ket}[1]{|#1\rangle}
\newcommand{\bra}[1]{\langle #1|}
\newcommand{\scalar}[2]{\langle #1|#2\rangle}
\newcommand{\proj}[1]{\left|#1\right\>\!\left\<#1\right|}
\newcommand{\norm}[1]{\|#1\|}
\newcommand{\ot}{\otimes}
\newcommand{\eps}{\epsilon}
\newcommand{\ra}{\rightarrow}

\renewcommand{\setminus}{\mathbin{-}}
\newcommand{\partitionof}{\mathbin{\vdash}}


\newcommand{\I}{\mathrm{I}}
\renewcommand{\H}{\mathrm{H}}
\renewcommand{\S}{\mathrm{S}}
\newcommand{\T}{\mathrm{T}}

%

\newcommand{\bI}{\mathbb{I}}
\newcommand{\bH}{\mathbb{H}}
\newcommand{\bS}{\mathbb{S}}
\newcommand{\bT}{\mathbb{T}}




%%%%%%%%%%%%%%%%%%%%%%%%%%%%%%%%%%%%%%%%%%%%%%%%%%%%%%%%%%%%

\title{Introduction to Quantum Algorithms Based on Group Representations}
\author{Pawel M. Wocjan\\[1ex]School of Electrical Engineering and Computer Science\\ University of Central Florida\\Orlando}
\email{wocjan@cs.ucf.edu}

\begin{document}

\maketitle

%%%%%%%%%%%%%%%%%%%%%%%%%%%%%%%%%%%%%%%%%%%%%%%%%%%%%%%%%%%%%%%%%%%%%%%%

\begin{slide}{Quantum Computing}

\end{slide}

%%%%%%%%%%%%%%%%%%%%%%%%%%%%%%%%%%%%%%%%%%%%%%%%%%%%%%%%%%%%%%%%%%%%%%%%

\begin{slide}{Black-box Functions}

Let $R$ and $S$ be two sets.

\medskip
A function $f : R \rightarrow S$ is called a black-box function if

\begin{itemize}
\item we don't know how $f$ is implemented 
\item the only way we can learn something about $f$ is to query $f$, i.e., to choose $r\in R$ and to obtain $f(r)$ (possibly multiple times; the queries could be adaptive)
\end{itemize}

Usually, the task is to determine some property of $f$.

\end{slide}

%%%%%%%%%%%%%%%%%%%%%%%%%%%%%%%%%%%%%%%%%%%%%%%%%%%%%%%%%%%%%%%%%%%%%%%%

\begin{slide}{Black-box Functions on Groups}

Let $G$ be a group and $S$ some set.

\bigskip
We say that the black-box function $f : G \rightarrow S$ hides a subgroup $H\le G$ if $f$ satisfies the following two conditions:

\begin{itemize}
\item $f$ is constant on the left cosets of $H$ in $G$, i.e.,
\[
f(g)=f(h) \mbox{ iff $gH=hH$}
\]
\item $f$ is different on different cosets, i.e., 
\[
f(g)\neq f(h) \mbox{ iff $gH\neq hH$}
\]
\end{itemize}

\end{slide}

%%%%%%%%%%%%%%%%%%%%%%%%%%%%%%%%%%%%%%%%%%%%%%%%%%%%%%%%%%%%%%%%%%%%%%%%

\begin{slide}{The Stabilizer Problem as a HSP}

\end{slide}

%%%%%%%%%%%%%%%%%%%%%%%%%%%%%%%%%%%%%%%%%%%%%%%%%%%%%%%%%%%%%%%%%%%%%%%%

\begin{slide}{Important Applications of HSP}

\end{slide}

%%%%%%%%%%%%%%%%%%%%%%%%%%%%%%%%%%%%%%%%%%%%%%%%%%%%%%%%%%%%%%%%%%%%%%%%

\begin{slide}{Complexity of the Hidden Subgroup Problem}

The task is to determine a generating set of the hidden subgroup $H$ by only querying $f$.

\bigskip
Classically, we need $\Omega(|G|)$ queries to solve the HSP.  

\bigskip
In contrast, we can solve the problem with high probability and running time $O(\log^c(|G|))$ on quantum computer for the following families of groups:

\begin{itemize}
\item
\item
\end{itemize}
\end{slide}

%%%%%%%%%%%%%%%%%%%%%%%%%%%%%%%%%%%%%%%%%%%%%%%%%%%%%%%%%%%%%%%%%%%%%%%%

\begin{slide}{Intro to Quantum Computing}

\end{slide}

%%%%%%%%%%%%%%%%%%%%%%%%%%%%%%%%%%%%%%%%%%%%%%%%%%%%%%%%%%%%%%%%%%%%%%%%

\begin{slide}{Pontrajagin-duality}
\textcolor{blue}{Theorem (Pontrajagin-duality)}

For every finite abelian group $G$ there is an isomorphism $\phi$ with $G\cong_{\phi} \hat{G}:={\rm Hom}(G,\C^\times)$.

\medskip
Let $\beta$ be the pairing between $G$ and $\hat{G}$ defined by 
\[
\beta : \left\{
\begin{array}{ccc}
G \times \hat{G} & \rightarrow & \C^\times \\
(g,\chi)         & \mapsto     & \chi(g)
\end{array}
\right.
\]

\end{slide}

%%%%%%%%%%%%%%%%%%%%%%%%%%%%%%%%%%%%%%%%%%%%%%%%%%%%%%%%%%%%%%%%%%%%%%%

\begin{slide}{Orthogonal complement of subgroups}

\medskip
We say that $g\in G$ is orthogonal to $\chi\in\hat{G}$, denoted by $g \perp \chi$, if $\beta(g,\chi)=1$, that is, $g\in{\rm ker}(\chi)$.

\medskip
We identify $\hat{A}$ and $A$ with respect to the isomorphism $\phi$ and define for a subgroup $H\le G$ an orthogoanl complement by
\[
H^{\perp} := \{y\in G \mid \forall x\in H : y \perp x\} 
\]
\end{slide}

%%%%%%%%%%%%%%%%%%%%%%%%%%%%%%%%%%%%%%%%%%%%%%%%%%%%%%%%%%%%%%%%%%%%%%%%

\begin{slide}{Fourier transform}
The Fouriertransform ${\rm DFT}_G$ over the abelian group $G$ is defined by
\[
{\rm DFT}_G := \frac{1}{\sqrt{|G|}} \sum_{x,y\in G} \beta(x,y) |y\>\<x|
\]

Lemma: For every subgroup $H\le G$ and $g\in G$ we have
\[
{\rm DFT}_G \frac{1}{\sqrt{H}} \sum_{x\in H} |g+x\> = \frac{1}{\sqrt{|H^\perp|}} \sum_{y\in H^\perp} \beta(g,y) |y\>
\]
\end{slide}

%%%%%%%%%%%%%%%%%%%%%%%%%%%%%%%%%%%%%%%%%%%%%%%%%%%%%%%%%%%%%%%%%%%%%%%%

\end{document}

